% Source: https://artofproblemsolving.com/community/c3h2986932_some_casework
% Stable identity: aops-topic-2986932-q01
A triple with two equal values $u,u,v$ ($u\ne v$) is never an arithmetic progression, while $u,u,u$ is. So a tuple is perfect iff (i) no value occurs three or more times, and (ii) the set of distinct values contains no three-term arithmetic progression.

By (i), at least $3$ distinct values occur. Subsets of $\{1,\dots,5\}$ with no $3$-term progression: the full set fails; among $4$-subsets only $\{1,2,4,5\}$ works; among the $10$ three-element subsets, the progressions $\{1,2,3\},\{2,3,4\},\{3,4,5\},\{1,3,5\}$ fail, leaving $6$.

\emph{Three values:} counts $2,2,1$: choose the single value ($3$ ways) and arrange: $\dfrac{5!}{2!2!}=30$. Total $6\cdot3\cdot30=540$.
\emph{Four values} $\{1,2,4,5\}$: counts $2,1,1,1$: choose the double ($4$) and arrange: $\dfrac{5!}{2!}=60$. Total $240$.

Answer: $540+240=\boxed{780}$.
